% GENERATED FILE. Edit canonical lessons, then run npm run generate:books.
% canonical-source-hash: fnv1a64:a3db12cb893e69e0
% canonical-lessons: RU-C319-legenda, RU-C319-mif, RU-C319-geroy, RU-C319-rasskazchik, RU-C319-sochinyat

\chapter{Legends and heroes}
\label{ch:ru-legends-and-heroes}
\hlchaptermodality{319}

\begin{chapteropening}
\textbf{By the end of this chapter:} \emph{I can talk about a legend, a myth, a hero and a storyteller, and say that someone made up a story.}

Complete the last lesson of chapter 319: I can talk about a legend, a myth, a hero and a storyteller, and say that someone made up a story.
\end{chapteropening}

\section[\texorpdfstring{\ru{легенда} (legénda)}{легенда (legénda)}]{\ru{легенда} (legénda) — a legend}

\label{lesson:RU-C319-legenda}



\begin{quote}
Before the new one: say the Russian for a ball, then the Russian for festivities.
\end{quote}

\subsection*{You'll want to know: \ru{легенда}}
\textbf{\ru{легенда}} (\emph{legénda}) — \textquotedblleft{}a legend\textquotedblright{}.

\begin{quote}
\ru{Об} \ru{этом} \ru{озере} \ru{есть} \ru{красивая} \ru{легенда}.
\end{quote}

\begin{quote}
\emph{Ob étom ózere yest krasívaya legénda.}
\end{quote}

\begin{quote}
There is a beautiful legend about this lake.
\end{quote}

Feminine. \emph{\ru{По} \ru{легенде}} — according to legend.

\begin{cousinweb}
\emph{\ru{легенда}} goes back to Medieval Latin \emph{legenda} \textquotedblleft{}things to be read\textquotedblright{}, from \emph{legere} \textquotedblleft{}to gather, to read\textquotedblright{} (Lewis \& Short); EtymDB lists no parent for the Russian word.
\end{cousinweb}

\begin{grammarlens}[title={What you've built}]
One more word for telling, explaining and describing at B1.
\end{grammarlens}

\subsection*{Your turn}
\begin{itemize}
\raggedright
  \item \emph{Say it:} \emph{legénda}
  \item \emph{Say it:} \emph{legénda}, once more
  \item \emph{Say it:} a sentence of your own with \emph{legénda}
  \item \emph{Recall:} say the Russian for a ball, then the Russian for festivities, then say \emph{legénda} again
\end{itemize}

\begin{tcolorbox}[breakable,colback=teal!4,colframe=teal!35!black,title={Before you move on}]
What does \ru{легенда} mean? (\textquotedblleft{}A legend\textquotedblright{}.) Say it once more.
\end{tcolorbox}
\section[\texorpdfstring{\ru{миф} (mif)}{миф (mif)}]{\ru{миф} (mif) — a myth}

\label{lesson:RU-C319-mif}



\begin{quote}
Before the new one: say the Russian for festivities, then the Russian for a legend.
\end{quote}

\subsection*{You'll want to know: \ru{миф}}
\textbf{\ru{миф}} (\emph{mif}) — \textquotedblleft{}a myth\textquotedblright{}.

\begin{quote}
\ru{Это} \ru{только} \ru{миф}.
\end{quote}

\begin{quote}
\emph{Éto tólko mif.}
\end{quote}

\begin{quote}
That is only a myth.
\end{quote}

Masculine. \emph{\ru{Мифы} \ru{Древней} \ru{Греции}} — the myths of ancient Greece.

\begin{cousinweb}
\emph{\ru{миф}} goes back to Greek \emph{mûthos} \textquotedblleft{}a word, a tale\textquotedblright{} (EtymDB).
\end{cousinweb}

\begin{grammarlens}[title={What you've built}]
One more word for telling, explaining and describing at B1.
\end{grammarlens}

\subsection*{Your turn}
\begin{itemize}
\raggedright
  \item \emph{Say it:} \emph{mif}
  \item \emph{Say it:} \emph{mif}, once more
  \item \emph{Say it:} a sentence of your own with \emph{mif}
  \item \emph{Recall:} say the Russian for festivities, then the Russian for a legend, then say \emph{mif} again
\end{itemize}

\begin{tcolorbox}[breakable,colback=teal!4,colframe=teal!35!black,title={Before you move on}]
What does \ru{миф} mean? (\textquotedblleft{}A myth\textquotedblright{}.) Say it once more.
\end{tcolorbox}
\section[\texorpdfstring{\ru{герой} (geróy)}{герой (geróy)}]{\ru{герой} (geróy) — a hero; a main character}

\label{lesson:RU-C319-geroy}



\begin{quote}
Before the new one: say the Russian for a legend, then the Russian for a myth.
\end{quote}

\subsection*{You'll want to know: \ru{герой}}
\textbf{\ru{герой}} (\emph{geróy}) — \textquotedblleft{}a hero; a main character\textquotedblright{}.

\begin{quote}
\ru{Главный} \ru{герой} \ru{книги} — \ru{маленький} \ru{мальчик}.
\end{quote}

\begin{quote}
\emph{Glávnyy geróy knígi — málenkiy málchik.}
\end{quote}

\begin{quote}
The main character of the book is a little boy.
\end{quote}

\emph{\ru{Главный} \ru{герой}} — the main character.

\begin{cousinweb}
\emph{\ru{герой}} came through French \emph{héros} and Latin \emph{hērōs} from Greek \emph{h\'{\={e}}rōs} \textquotedblleft{}a hero, a demigod\textquotedblright{} (EtymDB).
\end{cousinweb}

\begin{grammarlens}[title={What you've built}]
One more word for telling, explaining and describing at B1.
\end{grammarlens}

\subsection*{Your turn}
\begin{itemize}
\raggedright
  \item \emph{Say it:} \emph{geróy}
  \item \emph{Say it:} \emph{geróy}, once more
  \item \emph{Say it:} a sentence of your own with \emph{geróy}
  \item \emph{Recall:} say the Russian for a legend, then the Russian for a myth, then say \emph{geróy} again
\end{itemize}

\begin{tcolorbox}[breakable,colback=teal!4,colframe=teal!35!black,title={Before you move on}]
What does \ru{герой} mean? (\textquotedblleft{}A hero; a main character\textquotedblright{}.) Say it once more.
\end{tcolorbox}
\section[\texorpdfstring{\ru{рассказчик} (rasskázchik)}{рассказчик (rasskázchik)}]{\ru{рассказчик} (rasskázchik) — a storyteller, a narrator}

\label{lesson:RU-C319-rasskazchik}



\begin{quote}
Before the new one: say the Russian for a myth, then the Russian for a hero.
\end{quote}

\subsection*{You'll want to know: \ru{рассказчик}}
\textbf{\ru{рассказчик}} (\emph{rasskázchik}) — \textquotedblleft{}a storyteller, a narrator\textquotedblright{}.

\begin{quote}
\ru{Дедушка} \ru{был} \ru{отличным} \ru{рассказчиком}.
\end{quote}

\begin{quote}
\emph{Dédushka byl otlíchnym rasskázchikom.}
\end{quote}

\begin{quote}
Grandpa was an excellent storyteller.
\end{quote}

A woman is \emph{\ru{рассказчица}}.

\begin{cousinweb}
\emph{\ru{рассказчик}} is built on \emph{\ru{рассказать}} \textquotedblleft{}to tell\textquotedblright{}, from \emph{\ru{рас}-} and \emph{\ru{казать}} \textquotedblleft{}to show, to tell\textquotedblright{}, inherited from Proto-Slavic \emph{kazati} (EtymDB), with the ending \emph{-\ru{чик}} for a person.
\end{cousinweb}

\begin{grammarlens}[title={What you've built}]
One more word for telling, explaining and describing at B1.
\end{grammarlens}

\subsection*{Your turn}
\begin{itemize}
\raggedright
  \item \emph{Say it:} \emph{rasskázchik}
  \item \emph{Say it:} \emph{rasskázchik}, once more
  \item \emph{Say it:} a sentence of your own with \emph{rasskázchik}
  \item \emph{Recall:} say the Russian for a myth, then the Russian for a hero, then say \emph{rasskázchik} again
\end{itemize}

\begin{tcolorbox}[breakable,colback=teal!4,colframe=teal!35!black,title={Before you move on}]
What does \ru{рассказчик} mean? (\textquotedblleft{}A storyteller, a narrator\textquotedblright{}.) Say it once more.
\end{tcolorbox}
\section[\texorpdfstring{\ru{сочинять} (sochinyát)}{сочинять (sochinyát)}]{\ru{сочинять} (sochinyát) — to make up, to compose}

\label{lesson:RU-C319-sochinyat}



\begin{quote}
Before the new one: say the Russian for a hero, then the Russian for a storyteller.
\end{quote}

\subsection*{You'll want to know: \ru{сочинять}}
\textbf{\ru{сочинять}} (\emph{sochinyát}) — \textquotedblleft{}to make up, to compose\textquotedblright{}.

\begin{quote}
\ru{Мама} \ru{сочиняла} \ru{сказки} \ru{для} \ru{нас}.
\end{quote}

\begin{quote}
\emph{Máma sochinyála skázki dlya nas.}
\end{quote}

\begin{quote}
Mum made up fairy tales for us.
\end{quote}

Imperfective; perfective \emph{\ru{сочинить}} (\emph{sochinít}). Also of music or verse; in speech also \textquotedblleft{}to tell tales\textquotedblright{}: \emph{\ru{не} \ru{сочиняй}!}

\begin{cousinweb}
\emph{\ru{сочинять}} is built on \emph{\ru{со}-} \textquotedblleft{}together\textquotedblright{} and \emph{\ru{чин}} \textquotedblleft{}order, rank\textquotedblright{}, inherited from Proto-Slavic \emph{činŭ} (EtymDB): to put together in order.
\end{cousinweb}

\begin{grammarlens}[title={What you've built}]
One more word for telling, explaining and describing at B1.
\end{grammarlens}

\subsection*{Your turn}
\begin{itemize}
\raggedright
  \item \emph{Say it:} \emph{sochinyát}
  \item \emph{Say it:} \emph{sochinyát}, once more
  \item \emph{Say it:} a sentence of your own with \emph{sochinyát}
  \item \emph{Recall:} say the Russian for a hero, then the Russian for a storyteller, then say \emph{sochinyát} again
\end{itemize}

\begin{tcolorbox}[breakable,colback=teal!4,colframe=teal!35!black,title={Before you move on}]
What does \ru{сочинять} mean? (\textquotedblleft{}To make up, to compose\textquotedblright{}.) Say it once more.
\end{tcolorbox}
